\section{مقدمه}

بیماری‌های قلبی عروقی همچنان یکی از مهم‌ترین علل مرگ‌ومیر در سطح جهان به شمار می‌روند.
تشخیص زودهنگام و برآورد صحیح ریسک ابتلا، نقش کلیدی در پیشگیری از عوارض شدید و بهبود
کیفیت زندگی بیماران دارد. در سال‌های اخیر، روش‌های مبتنی بر داده و تکنیک‌های یادگیری
ماشین به ابزارهایی بسیار مهم برای تحلیل داده‌های پزشکی و پشتیبانی از تصمیم‌گیری بالینی
تبدیل شده‌اند.

در این پروژه از تکنیک‌های داده‌کاوی بر روی مجموعه‌دادهٔ شناخته‌شدهٔ بیماری قلبی از مخزن
\lr{UCI Machine Learning Repository} استفاده می‌شود. هدف اصلی این است که بررسی شود چگونه
ویژگی‌های مختلف بیماران، مانند سن، جنسیت، نوع درد قفسهٔ سینه، سطح کلسترول، فشار خون
استراحتی و حداکثر ضربان قلب، می‌توانند برای پیش‌بینی وجود بیماری قلبی به کار گرفته شوند.

تمرکز این مطالعه بر الگوریتم طبقه‌بندی \lr{K-Nearest Neighbors (KNN)} است که یک روش
ساده اما قدرتمند مبتنی بر فاصله محسوب می‌شود. مراحل اصلی پروژه شامل پیش‌پردازش داده‌ها،
تحلیل اکتشافی داده‌ها (\lr{EDA})، پیاده‌سازی طبقه‌بند \lr{KNN} برای ویژگی‌های ترکیبی
عددی و دسته‌ای، و ارزیابی عملکرد مدل با استفاده از معیارهای استاندارد طبقه‌بندی است.
نتایج این مطالعه، توان بالقوهٔ الگوریتم \lr{KNN} در مسائل طبقه‌بندی پزشکی را نشان می‌دهد
و اهمیت پیش‌پردازش مناسب و تحلیل دقیق ویژگی‌ها را در فرایند داده‌کاوی برجسته می‌سازد.